% HD-F1-T24. Proposed by /root; reconstructed by /root/transport.
% GPT-6 Astra Pro; 22 November 2026 UTC.
\section[HD-F1-T24: exact Gaussian alternating coefficient models]{HD-F1-T24: split alternating constituents are excluded for every scalar lift, and the Gaussian residual coefficient has an explicit rational symmetric-group model}
\label{sec:f1-gaussian-alternating-coefficient}

All split self-conjugate ordinary alternating
constituents are excluded at the remaining natural
sixfold source, including over \(K=\mathbf Q(i)\)
and for every original scalar lift. Consequently
every remaining Gaussian alternating coefficient
has the following exact form, on its original
\(K\)-space:
\[
 F\rho(g)F^{-1}=c(g)\rho_\lambda(\sigma(g)),\qquad
 \lambda\ne\lambda',\quad
 c:G\twoheadrightarrow\mu_4,\quad
 \sigma:G\longrightarrow S_n,\quad n\ge7.
 \tag{F1T24.1}
\]
Here \(F\) is an actual \(K\)-linear intertwiner,
\(\rho_\lambda\) is the rational
Young-symmetrizer representation of \(S_n\),
and \((c,\sigma)\) is an injective homomorphism
from the original coefficient group. Its image is
exactly one of
\[
 \mu_4\times A_n,\qquad \mu_4\times S_n,
 \qquad
 \{(z,s)\in\mu_4\times S_n:z^2=\operatorname{sgn}(s)\}.
 \tag{F1T24.2}
\]
All correlations in these displayed groups are retained.
The statement does not assert that their remaining
nonlinear ordinary modules yield a dominant source.

\paragraph{1. Original restriction and the explicit odd-centralizer map.}
By T21 and T23, the Gaussian alternating case has
actual perfect layer \(L\simeq A_n\), \(n\ge7\),
and \(P=ZL\) of index at most two in the original
\(G\). The averaging and evaluation proof in
T22, paragraph2, shows that the restriction of the
original absolutely irreducible coefficient module
to \(L\) has either one irreducible constituent
or two odd-conjugate constituents, each with
multiplicity one. That proof uses only the scalar
action of \(Z\) and the quotient of order at most
two; its conclusion applies equally to Gaussian
scalars.

Consider a self-conjugate partition \(\lambda\)
and its rational symmetric-group module
\[
 U=S^\lambda_{\mathbf C},\qquad
 U|_{A_n}=W\oplus W',\qquad
 \dim W=\dim W'=d.
 \tag{F1T24.3}
\]
Every odd permutation exchanges \(W\) and \(W'\),
by the exact index-two proof already read and used
in T22. No identification with complex conjugation
is made in this proof.

For \(\ell\in A_n\), its symmetric-group
conjugacy class is nonsplit precisely when its
centralizer contains an odd permutation. This
criterion follows directly from the conjugation
map \(S_n/C_{S_n}(\ell)\) and its restriction
to \(A_n\): the latter is surjective exactly
when \(S_n=A_n C_{S_n}(\ell)\).
It also gives explicit odd centralizers. An even
cycle is itself an odd permutation centralizing
\(\ell\). For two cycles of the same odd
length, exchanging corresponding letters of those
cycles is odd and centralizes \(\ell\). If all
cycle lengths are odd and distinct, a commuting
permutation preserves each cycle and is a product
of its powers, so is even. Thus the split classes
are exactly the classes with distinct odd cycle
lengths.

In the nonsplit case choose one such odd
centralizer \(s\). The actual operator
\[
 J_\ell=\rho_\lambda(s)|_W:W\xrightarrow{\sim}W',
 \qquad
 J_\ell\rho_W(\ell)=\rho_{W'}(\ell)J_\ell
 \tag{F1T24.4}
\]
is an isomorphism of the restricted cyclic
representations. It commutes with every power of
\(\ell\), so supplies directly the equality of
all power characters and all eigenspace
multiplicities. In particular, for every complex
scalar \(z\),
\[
 \operatorname{rank}_U(z\rho_\lambda(\ell)-1)
       =2\operatorname{rank}_W(z\rho_W(\ell)-1).
 \tag{F1T24.5}
\]
This proves the needed comparison for arbitrary
scalar lifts, with an explicit intertwiner.

\paragraph{2. Nonsplit classes and their actual transposition commutator.}
Let \(\ell\ne1\) have a nonsplit class and
choose a transposition \(\tau\) not commuting
with it, as in T22. Then
\(h=[\ell,\tau]\) is a 3-cycle or double
transposition. The content and class-sum calculation
of T22, equations (F1T22.7)--(F1T22.12), applies
to every self-conjugate partition and gives
\[
             \operatorname{rank}_U(h-1)\ge4d/5.
 \tag{F1T24.6}
\]
Set \(B=z\rho_\lambda(\ell)\) and
\(T=\rho_\lambda(\tau)\), for any nonzero
scalar \(z\). The exact identity
\[
 \rho_\lambda(h)-1
       =((B-1)T-T(B-1))B^{-1}T^{-1}
 \tag{F1T24.7}
\]
and (F1T24.5) give
\[
 \operatorname{rank}_W(z\rho_W(\ell)-1)\ge d/5.
 \tag{F1T24.8}
\]
The same proof applies with \(W'\) in place of
\(W\).

\paragraph{3. Split classes and a commutator inside the original alternating group.}
Now let \(\ell\) belong to a split class. Since
its cycle lengths are distinct and odd and
\(n\ge7\), one cycle has length at least five:
the only distinct odd lengths below five are
one and three, whose sum is four. On such a cycle
choose four successive letters
\(a_j=\ell^j(a_0)\), \(0\le j\le3\), and
take the actual element
\(\tau=(a_0\ a_1\ a_2)\in A_n\).
Using rightmost-first permutation composition,
\[
 [\ell,\tau]
   =(a_1\ a_2\ a_3)(a_0\ a_2\ a_1)
   =(a_0\ a_3\ a_1).
 \tag{F1T24.9}
\]
Thus this actual commutator is a nonidentity
3-cycle. It has at least four fixed letters.
A transposition of two fixed letters supplies
the odd centralizer in (F1T24.4), now for \(h\).
Consequently its two constituent moving ranks
are equal and each is half the rank on \(U\).
The three-cycle bound of T22 gives
\[
 \operatorname{rank}_W(h-1)
      =\tfrac12\operatorname{rank}_U(h-1)\ge d/2.
 \tag{F1T24.10}
\]
Both \(\ell\) and \(\tau\) are in \(A_n\),
so this time the commutator identity is on the
original constituent \(W\) itself:
\[
 \rho_W(h)-1=
 ((z\rho_W(\ell)-1)\rho_W(\tau)
 -\rho_W(\tau)(z\rho_W(\ell)-1))
       (z\rho_W(\ell))^{-1}\rho_W(\tau)^{-1}.
 \tag{F1T24.11}
\]
It proves
\(\operatorname{rank}_W(z\rho_W(\ell)-1)
\ge d/4\), for every nonzero scalar \(z\).
Again the assertion holds for \(W'\) as well.
Together with (F1T24.8), every nonidentity
scalar lift of a nonidentity permutation moves
at least \(d/5\) on either constituent.
For \(\ell=1\), a nonidentity scalar moves
all \(d\) dimensions.

\paragraph{4. Both original restriction patterns are excluded.}
With one split constituent, an outside element
of \(G/P\) would exchange it with the
inequivalent other constituent, so \(G=P\).
The preceding bounds apply directly to every
original nonidentity matrix on its \(m=d\)
dimensional coefficient module.

With two constituents, \(m=2d\). Inside \(P\)
the two moving ranks add, giving at least
\(m/5\). Every outside matrix exchanges the
two summands and has the form
\((x,y)\mapsto(Ay,Bx)\) with both off-diagonal
maps invertible. Projection from any eigenspace
to the first summand is injective: a vector
\((0,y)\) in that eigenspace satisfies
\(Ay=0\), hence \(y=0\).
Thus its eigenspaces have dimension at most
\(d\), and its moving rank is at least
\(m/2\). These arguments include every original
scalar choice.

In either pattern, all original nonidentity
inertias therefore move at least \(m/5\).
The original source equations are still
\[
 2m+6=\sum_{j=1}^b\operatorname{rank}(g_j-1),
                 \qquad b\ge12,\quad m\ge27.
 \tag{F1T24.12}
\]
They imply \(m\le15\), a contradiction.
This excludes the entire split-constituent
ordinary alternating branch for all scalars.

\paragraph{5. The residual rational model is defined over the actual Gaussian field.}
By the index-two representation description, the
remaining restriction to \(L=A_n\) is the
absolutely irreducible restriction of
\(S^\lambda\), with \(\lambda\ne\lambda'\).
Choose its rational Young-symmetrizer matrices.
The space
\[
 \operatorname{Hom}_{K[L]}
       (V,K\otimes_{\mathbf Q}S^\lambda)
 \tag{F1T24.13}
\]
is defined by linear equations over \(K\) in
the entries of an intertwiner. After extension
to \(\mathbf C\) it has dimension one by
Schur's lemma. Thus it has a nonzero \(K\)-point
\(F\), and this map is invertible because its
complexification is a nonzero map between
isomorphic irreducibles. This constructs the
actual \(K\)-linear comparison in (F1T24.1).

Conjugation on \(L\) gives the original
homomorphism \(\sigma:G\to S_n\), whose
kernel is \(Z\) and whose image contains
\(A_n\). The two matrices
\(F\rho(g)F^{-1}\) and
\(\rho_\lambda(\sigma(g))\) implement
the same conjugation of \(\rho_\lambda(L)\).
Their quotient is therefore a scalar
\(c(g)\in K^*\), by the same absolutely
irreducible centralizer calculation. The
representation identities give
\(c(gh)=c(g)c(h)\); finite order gives
\(c(g)\in\mu_4\). This proves (F1T24.1).
Faithfulness of the original representation
makes \((c,\sigma)\) injective.

The character field is exactly \(K\), by the
actual center map (F1T22.3). All
\(\rho_\lambda\)-traces are rational
integers. If the image of \(c\) were
\(\{1\}\) or \(\{1,-1\}\), all original
traces would be rational. Consequently
\(c(G)=\mu_4\).
Since \(L\) is perfect, \(c(L)=1\), and
its permutation image is all \(A_n\).
Thus \(H=(c,\sigma)(G)\) contains
\(\{1\}\times A_n\).
The image \(\overline H\) in
\(\mu_4\times(S_n/A_n)=C_4\times C_2\)
projects onto \(C_4\). Its kernel for that
projection is either trivial or all \(C_2\).
In the second case \(\overline H=C_4\times C_2\).
In the first it is the graph of one of the two
homomorphisms \(C_4\to C_2\), namely the
trivial map or reduction modulo two.
Taking the full inverse image, using
\(\{1\}\times A_n\subset H\), gives
exactly the three groups (F1T24.2).

\paragraph{Exact source and geometric scope.}
The new maps are the explicit odd-centralizer
intertwiner (F1T24.4), the two commutator
comparisons, and the \(K\)-defined intertwiner
(F1T24.13) with its actual scalar and permutation
characters. The universal rational induction
from T22 is used only for the proved rank
comparison. The original connected cover and
coefficient local system, primitive rational
corner and its integral-cleared Jacobian image,
restricted theta form and Poincare biextension
are unchanged by these coefficient
identifications. Their parabolic pairing,
de Rham and Hodge maps, period derivative,
Picard and Hilbert maps, cycle and
Grothendieck classes and Adams operations
retain the identity comparison and multiplier
one on the actual source. No new detector
component is counted. The remaining task in
this alternating branch is now confined to
the explicit non-self-conjugate rational
symmetric-group modules in (F1T24.1), with
one of the three exact coefficient images
in (F1T24.2).
